% concatenated from main.tex
\documentclass[11pt]{article}

% --- Packages ---
\usepackage{amsmath, amsthm, amssymb, amsfonts, mathtools}
\usepackage{mathrsfs}
\usepackage{hyperref}
\usepackage{enumitem}
\usepackage{graphicx}
\usepackage{tikz}
\usepackage{bbm}
\usepackage{geometry}
\geometry{margin=1in}

% --- Theorem environments ---
\newtheorem{theorem}{Theorem}[section]
\newtheorem{lemma}[theorem]{Lemma}
\newtheorem{proposition}[theorem]{Proposition}
\newtheorem{corollary}[theorem]{Corollary}
\theoremstyle{definition}
\newtheorem{definition}[theorem]{Definition}
\newtheorem{remark}[theorem]{Remark}
\newtheorem{example}[theorem]{Example}

% --- Macros ---
\newcommand{\R}{\mathbb{R}}
\newcommand{\Z}{\mathbb{Z}}
\newcommand{\OO}{\mathbb{O}}
\newcommand{\HH}{\mathbb{H}}
\newcommand{\N}{\mathbb{N}}
\newcommand{\E}{\mathbb{E}}
\newcommand{\Prob}{\mathbb{P}}
\newcommand{\calL}{\mathcal{L}}
\newcommand{\Lfull}{L^{\mathrm{full}}}
\newcommand{\Lhalf}{L^{\mathrm{half}}}
\newcommand{\LLL}{\mathfrak{L}}
\newcommand{\eps}{\varepsilon}
\newcommand{\indicator}{\mathbbm{1}}

\title{\textbf{Convergence of Half-Space Last Passage Percolation Away from the Boundary to the Directed Landscape}}
\author{Xinyi Zhang}
\date{\today}

\begin{document}

\maketitle

\begin{abstract}
In this note, we prove convergence of the half-space exponential last passage percolation (LPP) model, away from the boundary, to the directed landscape. Our approach couples the half-space and full-space LPP models and constructs two barrier events based on the monotonicity of last passage paths. Combining this coupling with moderate deviation estimates for both models and the known convergence of full-space LPP to the directed landscape, we establish the desired convergence. 
\end{abstract}


\section{Introduction}
\label{sec:intro}

The Kardar--Parisi--Zhang (KPZ) universality class describes the large-scale behavior of a wide family of random growth and interacting particle systems in $1+1$ dimensions. Over the past two decades, remarkable progress has been made in understanding the full-space KPZ universality class. For geometric and exponential last passage percolation (LPP) models with i.i.d.\ weights, Johansson~\cite{Johansson_2000} and Baik--Deift--Johansson~\cite{baik1999distribution} proved that the properly centered and scaled passage time from $(0,0)$ to $(n,n)$ converges in distribution to the GUE Tracy--Widom law. Beyond one-point fluctuations, the Airy$_2$ process was identified as the scaling limit of multi-point fluctuations along the spatial direction by Pr{\"a}hofer and Spohn~\cite{prahofer2002scale}. Subsequent works of Corwin, Quastel, and Remenik~\cite{corwin2013continuum} extended this analysis to continuum statistics using determinantal structure, while Corwin and Hammond~\cite{corwin2014brownian} established the Brownian Gibbs property of the Airy line ensemble, whose top curve is the Airy$_2$ process. More recently, Dauvergne, Ortmann, and Virág~\cite{dauvergne2022directed} completed this picture by constructing the Airy sheet and the directed landscape, the conjectural universal scaling limits governing joint fluctuations in all space-time directions for full-space models within the KPZ class.


New phenomena emerge when boundary conditions are introduced. In the half-space LPP model, one considers up-right paths constrained to $\HH=\{(i,j)\in\Z^2: i\ge j\}$ with different diagonal weights representing the boundary interaction. Baik and Rains~\cite{baik2000limiting} first discovered the crossover fluctuation distributions interpolating between the GOE and GUE Tracy--Widom distributions, corresponding to different strengths of the boundary. Later, Baik, Barraquand, Corwin, and Suidan~\cite{Baik_2018} identified a two-dimensional crossover kernel that generalizes the the crossover distributions found in \cite{forrester1999correlations} and \cite{sasamoto2004fluctuations}. 

In this note we relate the full-space and half-space LPP models by showing that, when the boundary is non-attractive, the boundary has no effect in the scaling limit for points that lie sufficiently far inside the bulk. More precisely, we prove that if the observation point is at distance $n^{2/3+\delta}$ from the boundary for some fixed $\delta>0$, the half-space last passage time converges to the same limit as in the full-space model. We only present our proof for the exponential LPP model, but the same argument applies to the geometric LPP as the three key moderate deviation estimates used in the exponential case all have geometric analogues.

When the boundary is non-attractive (in the exponential case, when $\alpha \ge \tfrac{1}{2}$), the full-space and half-space models share the same shape function, which can be viewed as the law of large number limit of the last passage value. Thus, the cost of deviating from the characteristic direction beyond the transversal fluctuation scale, which is of order $n^{2/3}$, to travel along the boundary is not compensated at leading order. In contrast, when the boundary becomes attractive (in the exponential case, when $\alpha < \tfrac{1}{2}$), the shape function for the half-space model is strictly larger, as the last passage path could benefit from traveling along the boundary even at the cost of deviating. In this regime, the portion of the path adhering to the boundary exhibits Gaussian fluctuations.

\section{Model Definitions and Main Results}
\label{sec:model}
Let $\prec$ denote the partial order on $\Z^2$ given by
$(x_1,y_1)\prec(x_2,y_2)$ if $x_1\le x_2$ and $y_1\le y_2$.
For both full-space and half-space exponential LPP with i.i.d.\ $\mathrm{Exp}(1)$ weights and non-attractive boundary, the length of a typical maximal path from $(x_1,y_1)$ to $(x_2,y_2)$ is encoded in the following deterministic shape function 
\[
d(x_1,y_1;x_2,y_2)
\;=\;\big(\sqrt{x_2-x_1+1}+\sqrt{\,y_2-y_1+1\,}\big)^2,
\qquad (x_1,y_1)\prec(x_2,y_2).
\]



For any subset $\OO \subset \Z^2$, let $\Pi_{\OO}[(x_1,y_1) \rightarrow (x_2,y_2)]$ denote the set of up-right lattice paths from $(x_1, y_1)$ to $(x_2,y_2)$ contained in $\OO$.

\begin{definition}[Full-space LPP]
Let $\{w_{i,j}\}_{i,j \in \Z^2}$ be i.i.d.\ exponential random variables with rate~1.  
Define the last passage time from $(x_1,y_1)$ to $(x_2,y_2)$ as
\[
\Lfull(x_1,y_1;x_2,y_2) = \max_{\pi \in \Pi[(x_1,y_1)\to(x_2,y_2)]} \sum_{(i,j)\in\pi \setminus (x_2,y_2)} w_{i,j},
\]
if $(x_1, y_1) \prec (x_2,y_2)$ and zero otherwise.
\end{definition}

\begin{definition}[Half-space LPP]
Let $\{w_{i,j}: i,j \in \Z \text{ and } i \geq j\}$ be independently distributed exponential random variables with rate~1 if $i > j$ and with rate $\alpha$ if $i = j$. Define the half space $\HH = \{(i,j) \in \Z^2, i \geq j\}$ and the half-space last passage time from $(x_1,y_1)$ to $(x_2,y_2)$ as
\[
\Lhalf(x_1,y_1;x_2,y_2) = \max_{\substack{ \pi \subset \HH \\\pi \in \Pi[(x_1,y_1)\to(x_2,y_2)]}} \sum_{(i,j)\in\pi \setminus (x_2,y_2)} w_{i,j},
\]
if $(x_1, y_1) \prec (x_2,y_2)$ and zero otherwise.
\end{definition}

Our main result establishes convergence to the directed landscape for
half-space exponential last passage percolation $n^{2/3+\delta}$ away from the boundary. We do not recall the definition of the directed landscape here and refer
the reader to \cite{dauvergne2022directed}.

\begin{theorem}[Half-space exponential LPP converges to the directed landscape]
\label{thm_half_to_DL}
Let $\R_+^4 = \{(x,s;r,t) \in \R^4: s<t\}$. Fix any $\delta>0$ and define the scaled half-space exponential last passage
percolation away from the boundary $L^{\mathrm{half},\delta}_n$ as a random function
on $\R_+^4$:
\begin{equation}\label{eq:Lhalfk_def}
\begin{split}
L^{\mathrm{half},\delta}_n(x,s;y,t)
=&\,2^{-4/3}n^{-1/3}\Lhalf(\lfloor ns+2^{5/3}n^{2/3}x\rfloor
+ \lfloor n^{2/3+\delta} \rfloor,\lfloor ns\rfloor;\,
\lfloor nt+2^{5/3}n^{2/3}y\rfloor+\lfloor n^{2/3+\delta} \rfloor,\lfloor nt\rfloor)\\
&-2^{2/3}n^{2/3}(t-s)-2^{4/3}n^{1/3}(y-x).
\end{split}
\end{equation}
Then $L^{\mathrm{half},\delta}_n$ converges to the directed landscape $\LLL$ in
distribution uniformly over compact subsets of $\R_+^4$.
\end{theorem}

We also recall the convergence of the full-space model \cite[Theorem 12.1]{dauv-vir-lis}, which provides the starting point of our analysis.

\begin{theorem}\label{thm_ELPP_to_DL}
For the following scaled full-space exponential last passage percolation, viewed as a random function on $\R_+^4$,
\begin{equation}\label{eq:Lhalfk_def}
\begin{split}
L^{\mathrm{full}}_n(x,s;y,t)
=&\,2^{-4/3}n^{-1/3}\Lfull(\lfloor ns+2^{5/3}n^{2/3}x\rfloor
,\lfloor ns\rfloor;\,
\lfloor nt+2^{5/3}n^{2/3}y\rfloor,\lfloor nt\rfloor)\\
&-2^{2/3}n^{2/3}(t-s)-2^{4/3}n^{1/3}(y-x).
\end{split}
\end{equation}
we know that $\Lfull_n$ converges to the directed landscape $\LLL$ in distribution uniformly over compact subsets on $\R_+^4$.
\end{theorem}

\section{Coupling Between Full-Space and Half-Space LPP}
\label{sec:coupling}

We couple the full-space model and the half-space model in the following way.  
Let $(\Omega,\mathcal{F},\Prob)$ be a probability space supporting two independent families of random variables
\[
\{W_{i,j}\}_{(i,j)\in\Z^2}\quad\text{and}\quad \{U_{i,i}\}_{i\in\Z},
\]
where the $\{W_{i,j}\}$ are i.i.d.\ $\mathrm{Exp}(1)$ and the $\{U_{i,i}\}$ are i.i.d.\ $\mathrm{Exp}(\alpha)$ for some fixed $\alpha \geq \frac{1}{2}$, independent of $\{W_{i,j}\}$.  
Define the full-space last passage time $L^{\mathrm{full}}$ using the weights $\{W_{i,j}\}_{(i,j)\in\Z^2}$, and the half-space last passage time $L^{\mathrm{half}}$ using the weights $\{W_{i,j}\}_{i>j}$ together with $\{U_{i,i}\}_{i\in\Z}$.  

Observe if we define the scaled full-space exponential last passage percolation away from the boundary as the following:
\begin{equation}\label{eq:Lfullk_def}
\begin{split}
L^{\mathrm{full},\delta}_n(x,s;y,t)
=&2^{-4/3}n^{-1/3}\Lfull(\lfloor ns+2^{5/3}n^{2/3}x\rfloor +\lfloor n^{2/3+\delta} \rfloor,\lfloor ns\rfloor;\,
                 \lfloor nt+2^{5/3}n^{2/3}y\rfloor+\lfloor n^{2/3+\delta} \rfloor,\lfloor nt\rfloor)\\
                 &
          -2^{2/3}n^{2/3}(t-s)-2^{4/3}n^{1/3}(y-x),
\end{split}
\end{equation}
then by Theorem~\ref{thm_ELPP_to_DL} and the translation invariance in distribution of exponential LPP, we have that $L^{\mathrm{full},\delta}$ converges to directed landscape in distribution uniformly over compact subsets on $\R_+^4$. By Theorem \ref{thm_ELPP_to_DL}, in order to establish the desired convergence for the half-space model, it suffices to prove the following coupling proposition:

\begin{proposition}\label{prop:diff}
For every compact $D\subset\R_+^4$,
\begin{equation}\label{eq:key_prob}
\lim_{n \rightarrow \infty}\Prob\!\left(L^{\mathrm{half},\delta}_n|_D \neq L^{\mathrm{full},\delta}_n|_D
  \right)= 0.
\end{equation}
\end{proposition}

To show that the probability in equation (\ref{eq:key_prob}) vanishes, we will construct two barrier events whose probability upper bound the probability in equation (\ref{eq:key_prob}) yet still converges to zero. To do so, we need the following lemma on the monotonicity of the rightmost last passage paths, whose continuous version can be found in \cite[Lemma 3.6]{dauvergne2022directed}.

\begin{lemma}\label{lemma_monotonicity}
    Let $\OO = \HH$ or $\Z^2$. Let $(x_1,y),(w_1,z),(x_2,y),(w_2,z) \in \OO$ such that $(x_1,y) \prec (w_1,z)$, $(x_2,y) \prec (w_2,z)$, and $x_1 \leq x_2, w_1 \leq w_2$. Let $P_{\OO}[(x_1,y) \rightarrow (w_1,z)] \subset \Pi_{\OO}[(x_1,y) \rightarrow (w_1,z)]$ be the set of path $\pi$ that maximizes the last passage value in $\OO$, i.e.,
    \[
    \sum_{(i,j) \in \pi} w_{i,j} = \max_{\substack{ \sigma \subset \OO \\\sigma \in \Pi[(x_1,y)\to(w_1,z)]}} \sum_{(i,j)\in \sigma \setminus (w_1,z)} w_{i,j}.
    \]
    Let $\pi^+(x_1,y;w_1,z) \in P_{\OO}[(x_1,y) \rightarrow (w_1,z)]$ denote the rightmost maximizing path and define the leftmost $x$-coordinate in path $\pi$ on the level $j$
    \begin{equation}
        \text{Leftmost}_j \left[\pi \right] = \begin{cases}
i & \text{if } (i-1,j) \notin \pi \text{ and } (i,j) \in \pi\\
-\infty & \text{if } (i,j) \notin \pi \text{ for all } i \in \Z \\
\end{cases}
    \end{equation}
    Then for any $j \in \Z$,
    \[\text{Leftmost}_j \left[\pi^+(x_1,y;w_1,z) \right] \leq \text{Leftmost}_j \left[\pi^+(x_2,y;w_2,z) \right].\]
\end{lemma}


\begin{proof}
Write $\pi_1:=\pi^+(x_1,y;w_1,z)$ and $\pi_2:=\pi^+(x_2,y;w_2,z)$. 
Assume for contradiction that there exists $j$ with $\text{Leftmost}_j[\pi_1]>\text{Leftmost}_j[\pi_2]$. Planarity implies that there exists two points $(i_1,j_1), (i_2, j_2) \in \pi_1 \cap \pi_2, (i_1,j_1) \prec (i_2,j_2)$ such that both $\pi_1$ and $\pi_2$ pass through $(i_1,j_1)$ and $(i_2,j_2)$ and on the level $j \in [j_1 + 1, j_2]$, $\pi_1$ is strictly to the right of $\pi_2$. 

We truncate the two paths into three segments each. Let $\pi_1 =
\pi_1^{(1)} \cup \pi_1^{(2)} \cup \pi_1^{(3)}$ where

\begin{enumerate}[label=(\arabic*), itemsep=2pt, leftmargin=2em]
    \item $\pi_1^{(1)}$ starts at $(x_1,y)$ and ends at $(i_1,j_1)$;
    \item $\pi_1^{(2)}$ starts at $(i_1+1,j_1)$ and ends at $(i_2,j_2-1)$;
    \item $\pi_1^{(3)}$ starts at $(i_2,j_2)$ and ends at $(w_1,z)$.
\end{enumerate}

Similarly, write $\pi_2 =
\pi_2^{(1)} \cup \pi_2^{(2)} \cup \pi_2^{(3)}$ where
\begin{enumerate}[label=(\arabic*), itemsep=2pt, leftmargin=2em]
    \item $\pi_2^{(1)}$ starts at $(x_2,y)$ and ends at $(i_1,j_1)$;
    \item $\pi_2^{(2)}$ starts at $(i_1,j_1+1)$ and ends at $(i_2-1,j_2)$;
    \item $\pi_2^{(3)}$ starts at $(i_2,j_2)$ and ends at $(w_2,z)$.
\end{enumerate}

Because $\pi_1$ is a maximizing path, we must have
\begin{equation}\label{eq:1}
    \sum_{(i,j) \in \pi_1^{(2)}} w_{i,j} \geq \sum_{(i,j) \in \pi_2^{(2)}} w_{i,j}.
\end{equation}
This implies that $\pi_2^{(1)} \cup \pi_1^{(2)} \cup \pi_2^{(3)}$ is also an element in $P_{\OO}[(x_2,y) \rightarrow (w_2,z)]$ and thus contradicts the fact that $\pi_2$ is the rightmost maximizing path.
\end{proof}

For any compact set $D \in \R^4_+$, we know that $D \subset [-M,M]^4$ for some large $M \in \Z$. Let us fix $\ell \in (0,1)$ and consider $n$ large enough such that $n^{2/3}M < n^{2/3+\delta}(1-\ell)$. Let $(x_1, y_1) = (\lfloor \ell n^{2/3+\delta} \rfloor-Mn, -Mn)$ and $(x_2, y_2) = (\lfloor \ell n^{2/3+\delta} \rfloor+Mn, Mn)$. Let $\text{Diagonal} = \{(i,i): i \in \Z\}$ and 
\begin{equation}
\begin{aligned}
    A(\HH) &= \{\omega \in \Omega: \text{for every } \pi \in P_{\HH}[(x_1,y_1) \rightarrow (x_2,y_2)], \pi \cap \text{Diagonal} \neq \emptyset\}\\
    A(\Z^2) &= \{\omega \in \Omega: \text{for every } \pi \in P_{\Z^2}[(x_1,y_1) \rightarrow (x_2,y_2)], \pi \cap \text{Diagonal} \neq \emptyset\}.
\end{aligned} 
\end{equation}

On the complement of $A(\HH)\cup A(\Z^2)$, there exists a maximizing path
$\pi^\star \in P_{\Z^2}[(x_1,y_1)\to(x_2,y_2)]$ and a maximizing path $\pi_\star \in P_{\HH}[(x_1,y_1)\to(x_2,y_2)]$ which avoid the diagonal. By Lemma~\ref{lemma_monotonicity}, we see that all the rightmost maximizing paths between endpoints $(x,s)$ and $(y,t)$ such that $(x,s;y,t) \in D$ avoid the diagonal. By the construction of our coupling, the full-space and half-space models
use identical weights at all off-diagonal sites. Therefore, for all $(x,s;y,t) \in D$,
\[
L^{\mathrm{half}}(x,s;y,t)
= L^{\mathrm{full}}(x,s;y,t)
\]
 on the complement of $A(\HH)\cup A(\Z^2)$. This implies
\[
\Prob\!\left(L^{\mathrm{half},\delta}_n|_D \neq L^{\mathrm{full},\delta}_n|_D\right)
\le \Prob\!\left(A(\HH)\cup A(\Z^2)\right) \leq \Prob(A(\HH)) + \Prob(A(\Z^2)).
\]

\section{Moderate Deviation Estimates}
\label{sec:full_moderate}
We will state and derive a few moderate deviation results in this section in order to upper bound the probabilities $\Prob(A(\HH))$ and $\Prob(A(\Z^2))$.

The following theorem comes from \cite[Theorem 2]{Small_Deviation} and is rephrased in terms of exponential last passage time in \cite[Theorem 2.2]{Nonexistence}

\begin{proposition}[Full-space one-point moderate deviation]
\label{prop:moderate}
For any $K > 1$, there exists two strictly positive constants $C(K), c(K)$ such that for all $m,n \geq 1$ with $K^{-1} < \frac{m}{n} < K$ and all $r > 0$ we have:
\begin{equation}
    \Prob\left(\Lfull(0,0;m,n) - (\sqrt{m} + \sqrt{n})^2 \geq rn^{1/3}\right) \leq Ce^{-c \min\{r^{3/2},rn^{1/3}\}}
\end{equation}
\begin{equation}
    \Prob\left(\Lfull(0,0;m,n) - (\sqrt{m} + \sqrt{n})^2 \leq -rn^{1/3}\right) \leq Ce^{-c r^3}.
\end{equation}
\end{proposition}

The next result gives an upper bound on the lower tail of constrained last passage value. Let us first define a cylinder $C\left[(w,z), \xi,\gamma\right]$ of width $2\sqrt{2}\gamma n^{2/3}$ from $(w,z)$ to $(w+\xi n, z+\xi n)$ in $\Z^2$:
\begin{equation}\label{eq:cylinder_def}
C_{\gamma}\big[(w,z),\xi\big]
:=\Big\{(x,y)\in\Z^2:
\ 2z \le x+y \le 2z+2\xi n,\ 
z-w-2\gamma n^{2/3} \le -x+y \le z-w+2\gamma n^{2/3}
\Big\}.
\end{equation}
Moreover, let $L^{C_\gamma}(w,z;w+\xi n, z+ \xi n)$ denote the maximum weight over all path from $(w,z)$ to $(w+\xi n, z+ \xi n)$ constrained inside $C_\gamma\left[(w,z), \xi\right]$. 

Now we are ready to introduce the following proposition from \cite[Proposition 3.7]{basu2022interlacing}:

\begin{proposition}[Full-space constrained moderate deviation]
\label{prop:constrained_moderate}
Fix $L_1, L_2 > 0$. Let $L_1 \leq \gamma \leq L_2$. Then there exists positive constants $\theta_0, n_0$, which depend on $L_1, L_2$, and an absolute constant $c>0$ such that for $n > n_0$ and $\theta > \theta_0$,
\begin{equation}
    \Prob\left(L^{C_\gamma}(1,1;\xi n, \xi n) - 4 \xi n \leq - \theta n^{1/3}\right) \leq e^{-c \gamma \theta}.
\end{equation}
\end{proposition}

Lastly, we need to derive an upper bound on the upper tail of half-space exponential LPP. We will use the key observation in \cite{Pfaffian_Schur} that the CDF of the half-space exponential LPP can be written as a Fredholm Pfaffian.



Let us first introduce the definition of a Fredholm Pfaffian.
\begin{definition}
Let \( K(x,y) \) be a \(2\times 2\) matrix-valued kernel on a measure space \((\mathbb{X},\mu)\), written as
\[
K(x,y) =
\begin{pmatrix}
K_{11}(x,y) & K_{12}(x,y)\\[2pt]
K_{21}(x,y) & K_{22}(x,y)
\end{pmatrix}, \qquad x,y \in \mathbb{X}.
\]
Let \(J(x,y)\) be defined by
\[
J(x,y) = \mathbf{1}_{x=y}
\begin{pmatrix}
0 & 1\\[2pt]
-1 & 0
\end{pmatrix}.
\]
We say that \(K\) is skew-symmetric if, for every collection of points
\(x_1,\ldots,x_{2k} \in \mathbb{X}\), the block matrix
\((K(x_i,x_j))_{i,j=1}^{2k}\) is skew-symmetric.
The \emph{Fredholm Pfaffian} of \(K\) is defined, whenever the series converges, by
\[
\operatorname{Pf}(J + K)_{\mathbb{L}^2(\mathbb{X},\mu)}
= 1 + \sum_{k=1}^{\infty} \frac{1}{k!}
\int_{\mathbb{X}^k}
\operatorname{Pf}\!\big( K(x_i,x_j) \big)_{i,j=1}^{k}
\, d\mu(x_1)\cdots d\mu(x_k).
\]
For a finite \(2k\times 2k\) skew-symmetric matrix \(A = (a_{ij})\),
its Pfaffian is given by
\[
\operatorname{Pf}(A)
= \frac{1}{2^k k!}
\sum_{\sigma \in S_{2k}}
\mathrm{sgn}(\sigma)
\, a_{\sigma(1)\sigma(2)} a_{\sigma(3)\sigma(4)} \cdots a_{\sigma(2k-1)\sigma(2k)}.
\]
\end{definition}

Additionally, we need the following lemma \cite[Lemma 2.5]{Pfaffian_Schur} which is proved by Hadamard's inequality and the fact that $\operatorname{Pf}(A) = \sqrt{\det(A)}$ for any skew-symmetric matrix $A$.

\begin{lemma}\label{lemma_pfaffian}
Let \(K(x,y)\) be a \(2\times 2\) skew-symmetric matrix-valued kernel.
Suppose there exist constants \(C>0\) and real numbers \(a>b\ge 0\) such that
\[
|K_{11}(x,y)| \le C e^{-a(x+y)}, \qquad
|K_{12}(x,y)| \le C e^{-a x + b y}, \qquad
|K_{22}(x,y)| \le C e^{b(x+y)}.
\]
Then, for every integer \(k > 0\),
\[
\bigl|\,\operatorname{Pf}\!\big(K(x_i,x_j)\big)_{i,j=1}^k\,\bigr|
\le (2k)^{k/2} C^k \prod_{i=1}^{k} e^{-(a-b)x_i}.
\]
\end{lemma}

Finally, we are ready to state the key proposition \cite[Proposition 1.6]{Pfaffian_Schur} and \cite[Lemma 6.4]{Pfaffian_Schur}:

\begin{proposition}\label{prop:pfaffian_schur}
For all $\alpha \geq 1/2$, $r \in \R$ and every positive integer $m$, 
    \begin{equation}
        \Prob\left(\Lhalf(1,1;m,m) -4m < 2^{4/3}m^{1/3}r\right) = \operatorname{Pf}[J-K^{\text{exp}, m}]_{\mathbb{L}^2(r,\infty)}.
    \end{equation}
For $\alpha = 1/2$, the hypotheses of Lemma \ref{lemma_pfaffian} are satisfied for $m$ large enough.
\end{proposition}

The exact formulas $K^{\text{exp}, m}$ is slightly different for the case that $\alpha > \frac{1}{2}$ and the case that $\alpha=\frac{1}{2}$. Since we will not use the
explicit form of $K^{\mathrm{exp},m}$, we omit it for simplicity. We now use Lemma~\ref{lemma_pfaffian} to control the convergence of the Fredholm Pfaffian series and to derive a quantitative upper bound for the tail of $\Lhalf(1,1;m,m)$.

\begin{proposition}[Half-space upper tail estimate]\label{prop_upper_tail}
For all $\alpha \geq 1/2$, there exist constants $c, C > 0$ and $m_0 \in \N$ such that, for all $m \geq m_0$ and all $r > 0$,
\[
\Prob\big(\Lhalf(1,1;m,m) - 4m \ge m^{1/3} r\big)
\le C \, e^{-c r}.
\]
\end{proposition}

\begin{proof}
By Proposition~\ref{prop:pfaffian_schur}, we know that
\[
\Prob\big(\Lhalf(1,1;m,m) - 4m \geq 2^{4/3}m^{1/3} r\big)
= 1- \operatorname{Pf}[J - K^{\mathrm{exp},m}]_{\mathbb{L}^2(r,\infty)}.
\]
where this Pfaffian can be expanded as a convergent series:
\[
1 - \operatorname{Pf}[J - K^{\mathrm{exp},m}]_{\mathbb{L}^2(r,\infty)}
= -\sum_{k=1}^{\infty} \frac{(-1)^k}{k!}
\int_{(r,\infty)^k} 
\operatorname{Pf}\!\big(K^{\mathrm{exp},m}(x_i,x_j)\big)_{i,j=1}^{k}
\, dx_1 \cdots dx_k.
\]
Lemma~\ref{lemma_pfaffian} guarantees that for $\alpha = 1/2$, each term of this expansion satisfies the bound
\[
\Big|
\operatorname{Pf}\!\big(K^{\mathrm{exp},m}(x_i,x_j)\big)_{i,j=1}^{k}
\Big|
\le (2k)^{k/2} C^k \prod_{i=1}^{k} e^{-(a-b)x_i},
\]
where the constants $a > b \ge 0$ are independent of $m$. Integrating this bound term by term yields
\[
\Big|
\int_{(r,\infty)^k} 
\operatorname{Pf}\!\big(K^{\mathrm{exp},m}(x_i,x_j)\big)_{i,j=1}^{k}
\, dx_1 \cdots dx_k
\Big|
\le (2k)^{k/2} C^k (a-b)^{-k} e^{-k(a-b)r}.
\]
Since the resulting series is absolutely summable, we obtain
\[
1 - \operatorname{Pf}[J - K^{\mathrm{exp},m}]_{\mathbb{L}^2(r,\infty)}
\le \sum_{k=1}^{\infty} \frac{1}{k!} (2k)^{k/2} C^k e^{-k(a-b)r}
\le C'e^{-(a-b)r},
\]
for some $C'>0$. Lastly, since exponential random variables are stochastically ordered in their rate
parameter, $\Lhalf(1,1;m,m)$ with $\alpha > 1/2$ is stochastically dominated by $\Lhalf(1,1;m,m)$ with $\alpha = 1/2$. Thus, the same upper tail applies.
\end{proof}

\section{Convergence Away from the Boundary}
\label{sec:mainproof}

Recall the endpoints
\[
(x_1,y_1)=\big(\lfloor \ell n^{2/3+\delta}\rfloor-Mn,\,-Mn\big),\qquad
(x_2,y_2)=\big(\lfloor \ell n^{2/3+\delta}\rfloor+Mn,\,Mn\big),
\]
with fixed $M\in\mathbb Z_{>0}$, fixed $\ell\in(0,1)$, and $n$ large
enough that $n^{2/3}M<n^{2/3+\delta}(1-\ell)$.


We now prove the following inequality for the shape function, showing that deviations beyond the typical transversal scale $n^{2/3}$ produce a gap between the free energy at leading order.
\begin{lemma}\label{lem:diag-deficit}
Let $I = [ \lfloor \ell n^{2/3+\delta} \rfloor -Mn,Mn].$ There exist positive constants $n_0, c_0, c_1$ such that the following holds for all $n \geq n_0$.

\textbf{(i)}
For any $i,j\in I$ satisfying $i\le j\le i+ n^{1/3}$, we have
\begin{equation}
d(x_1,y_1;j,j)+d(i,i;x_2,y_2)
\le 8Mn - c_0 n^{1/3+2\delta}.
\end{equation}

\textbf{(ii)}
For any $i,j,s,t\in I$ satisfying $i \leq j\le s\le t$, $t \leq s+ n^{1/3}$, and $j \leq i+ n^{1/3}$, we have
\begin{equation}\label{eq:hs}
d(x_1,y_1;j,j)+d(i,i;t,t)+d(s,s;x_2,y_2)
\le 8Mn - c_1 n^{1/3+2\delta}.
\end{equation}
\end{lemma}

\begin{proof}
    One can easily check that for $A > \max\{0,B\}$, 
    \[
    \sqrt{A^2 - B^2} \leq A - \frac{B^2}{2A}.
    \]
    Then,
    \[
    \begin{aligned}
    d(x_1,y_1;j,j) &= (\sqrt{j-x_1+1} + \sqrt{j-y_1+1})^2\\
    &= 2j - x_1 -y_1+2 + \sqrt{(2j-x_1-y_1+2)^2 - (y_1-x_1)^2}\\
    &\leq 2(2j-x_1-y_1+2) -\frac{\lfloor \ell n^{2/3 + \delta} \rfloor^2}{2(2j-x_1-y_1+2)}.
    \end{aligned}
    \]
    Similarly, 
    \[
    \begin{aligned}
    d(i,i;x_2,y_2) \leq 2(x_2+y_2 - 2i+2) -\frac{\lfloor \ell n^{2/3 + \delta} \rfloor^2}{2(x_2+y_2 - 2i+2)}.
    \end{aligned}
    \]
    Therefore,
    \[
    \begin{aligned}
    d(x_1,y_1;j,j) +d(i,i;x_2,y_2)  &\leq  8Mn +4j -4i -\frac{\lfloor \ell n^{2/3 + \delta} \rfloor^2}{4Mn+2}\\
    &\leq 8Mn + 4n^{1/3} - c_0 n^{1/3+2\delta}.
    \end{aligned}
    \]
    We absorb the term $4n^{1/3}$ in to the constant $c_0$ by choosing $n_0$ large enough. The proof for \eqref{eq:hs} follows analogously.
\end{proof}

\begin{proof}[Proof of Proposition~\ref{prop:diff}]
If $\omega \in A(\Z^2)$, then there exists a maximal path $\pi(\omega) \in P_{\Z^2}[(x_1,y_1) \rightarrow (x_2,y_2)]$ such that $(i,i) \in \pi(\omega)$ for some $i \in [\lfloor \ell n^{2/3+\delta} \rfloor -Mn,Mn].$ By the definition of last passage path, we know that
\begin{equation}
    \Lfull(x_1,y_1; x_2,y_2) \leq \Lfull(x_1,y_1; i,i) + \Lfull(i,i;x_2,y_2).
\end{equation}
Let us consider $n$ large enough so that $n^{1/3} > m_0$ for the constant $m_0$ in Proposition \ref{prop_upper_tail}. Let $w_0 < w_1 < \cdots < w_{k(n)}$ be a sequence of integers such that $w_0=\lfloor \ell n^{2/3+\delta} \rfloor -Mn, w_1 = \lfloor \ell n^{2/3+\delta} \rfloor -Mn + \lfloor n^{1/3} \rfloor , w_{k(n)-1} = Mn -\lfloor n^{1/3} \rfloor$ and $w_{k(n)} = Mn$. Moreover, $m_0 \leq w_{j+1} - w_j \leq n^{1/3}$ for all $0 \leq j \leq k(n)-1$.

If $w_j \leq  i < w_{j+1}$, then
\begin{equation}
    \Lfull(x_1,y_1; i,i) + \Lfull(i,i;x_2,y_2) \leq \Lfull(x_1,y_1; w_{j+1},w_{j+1}) + \Lfull(w_j,w_j;x_2,y_2).
\end{equation}

Now we are ready to bound $\Prob(A(\Z^2))$. Recall that $L^{C_\gamma}(w,z;w+\xi n, z+ \xi n)$ denote the maximum weight over all path from $(w,z)$ to $(w+\xi n, z+ \xi n)$ constrained inside $C_\gamma\left[(w,z), \xi\right]$. Choose $\gamma > \theta_0$ for the constant $\theta_0$ in Proposition \ref{prop:constrained_moderate} and $n$ large enough such that $8\gamma n^{2/3} < \ell n^{2/3+\delta}$.
\[
    \begin{aligned}
        \Prob(A(\Z^2)) &\leq \Prob\left(L^{C_\gamma}(x_1,y_1; x_2,y_2) < \max_{0 \leq j \leq k(n)-1} \left\{\Lfull(x_1,y_1; w_{j+1},w_{j+1}) + \Lfull(w_j,w_j;x_2,y_2)\right\} \right)\\
        &\leq \Prob\left(L^{C_\gamma}(x_1,y_1; x_2,y_2) - 8Mn < -\frac{1}{2}c_0 n^{1/3+2\delta} \right)\\ 
         &\quad + \Prob\left( \max_{0 \leq j \leq k(n)-1}\left\{\Lfull(x_1,y_1; w_{j+1},w_{j+1}) + \Lfull(w_j,w_j;x_2,y_2)\right\} -8Mn > -\frac{1}{2}c_0 n^{1/3+2\delta} \right)\\
         &\leq \Prob\left(L^{C_\gamma}(x_1,y_1; x_2,y_2) - 8Mn < -\frac{1}{2}c_0 n^{1/3+2\delta} \right)\\ 
         &\quad +  2Mn \max_{0 \leq j \leq k(n)-1}\Prob\left(\Lfull(x_1,y_1; w_{j+1},w_{j+1}) + \Lfull(w_j,w_j;x_2,y_2) -8Mn > -\frac{1}{2}c_0 n^{1/3+2\delta} \right)
    \end{aligned}
\]

By Proposition \ref{prop:constrained_moderate}, we know that
\[
\Prob\left(L^{C_\gamma}(x_1,y_1; x_2,y_2) - 8Mn < -\frac{1}{2}c_0 n^{1/3+2\delta} \right) \leq e^{-cn^{2\delta}}
\]
for some constant $c$ and all $n$ large enough.

On the other hand, by Lemma \ref{lem:diag-deficit}, we have 
\begin{equation}\label{eq:full}
\begin{split}
    &\Prob\left(\Lfull(x_1,y_1; w_{j+1},w_{j+1}) + \Lfull(w_j,w_j;x_2,y_2) -8Mn > -\frac{1}{2}c_0 n^{1/3+2\delta} \right)\\
    &\leq \Prob\left(\Lfull(x_1,y_1; w_{j+1},w_{j+1}) - d(x_1,y_1;w_{j+1},w_{j+1})+ \Lfull(w_j,w_j;x_2,y_2) - d(w_j,w_j;x_2,y_2)  >  \frac{1}{2}c_0 n^{1/3+2\delta} \right)\\
    &\leq \Prob\left(\Lfull(x_1,y_1; w_{j+1},w_{j+1}) - d(x_1,y_1;w_{j+1},w_{j+1}) > \frac{1}{4}c_0 n^{1/3+2\delta}\right) \\
    & \quad + \Prob\left(\Lfull(w_j,w_j;x_2,y_2) - d(w_j,w_j;x_2,y_2)  >  \frac{1}{4}c_0 n^{1/3+2\delta} \right).
\end{split}
\end{equation}
Let $K > 1$ to be the constant in Proposition \ref{prop:moderate} that will be chosen later. If $K^{-1} \leq \frac{y_2 - w_j}{x_2- w_j}, \frac{w_{j+1}-x_1}{w_{j+1}-y_1} \leq K$, then we can directly apply Proposition \ref{prop:moderate} and get an upper bound of $Ce^{-c \min \{n^{3\delta}, n^{1/3 + 2\delta}\}}.$ If not, we apply the following cheap bound
\[
\begin{aligned}
    &\Prob\left(\Lfull(1,1;m,m') - d(1,1;m,m') > 0 \right)\\
    &\leq \binom{m'+m-2}{m-1}\Prob\left( \sum_{i=1}^{m+m'} W_i - m-m' >  2\sqrt{mm'} \right) \\
    &\leq \exp\left( \log \binom{m'+m-2}{m-1}-\sqrt{mm'}\right)\\
\end{aligned}
\]
where $W_i$ are i.i.d. Exp(1) random variables. We use the fact that $\E[e^{W_1/2}] = 2$ and exponential Markov inequality in the last inequality. In our case, we have $m,m' \geq n^{1/3}$ and $0< \frac{m'}{m} < K^{-1}$. By Stirling's approximation,
\[
\log \binom{m+m'-2}{m-1} \leq m\left[\left(1+\frac{m'}{m}\right)\log\left(1+\frac{m'}{m}\right) - \frac{m'}{m}\log\left(\frac{m'}{m}\right) \right] + \mathcal{O}(\log m).
\]
Thus, we can choose $K^{-1}$ such that $\log \binom{m'+m-2}{m-1}-\sqrt{mm'} \leq -\sqrt{mm'}/2 \leq - n^{1/3}/2$. Thus, \eqref{eq:full} is upper bounded by $Ce^{-c {\min \{n^{2\delta}, n^{1/3} \}}}$.

Now we will prove the analogous statement for the half-space case. Let 
\[
\Lhalf|_{\text{full}}(x_1,y_1;i,i) = \max_{\substack{ \pi \subset \HH; \pi \cap \text{Diagonal} = (i,i)\\\pi \in \Pi[(x_1,y_1)\to(i,i)]}} \sum_{(u,v)\in\pi \setminus (i,i)} w_{u,v}
\]
\[
\Lhalf|_{\text{full}}(i,i;x_2,y_2) = \max_{\substack{ \pi \subset \HH; \pi \cap \text{Diagonal} = (i,i)\\\pi \in \Pi[(i,i)\to(x_2,y_2)]}} \sum_{(u,v)\in\pi \setminus (i,i),(x_2,y_2)} w_{u,v}
\]

Let
\[
L^R_{i,j} = \Lhalf|_{\text{full}}(x_1,y_1; w_{i+1},w_{i+1}) + \Lhalf(w_i,w_i;w_{j+1},w_{j+1}) + \Lhalf|_{\text{full}}(w_j,w_j;x_2,y_2),
\]
\[
L_{i,j} = \Lfull(x_1,y_1; w_{i+1},w_{i+1}) + \Lhalf(w_i,w_i;w_{j+1},w_{j+1}) + \Lfull(w_j,w_j;x_2,y_2).
\]
Since $L_{i,j} \geq L^R_{i,j}$ for all $0 \leq i\leq j \leq k(n)-1$, we can now bound $\Prob(A(\HH))$ by the following 
\begin{equation}
    \begin{aligned}
        &\Prob\left(L^{C_\gamma}(x_1,y_1; x_2,y_2) < \max_{0 \leq i \leq  j \leq k(n)-1} L^R_{i,j} \right)\\
        &\leq \Prob\left(L^{C_\gamma}(x_1,y_1; x_2,y_2) < \max_{0 \leq i \leq j \leq k(n)-1} L_{i,j}\right)\\
         &\leq \Prob\left(L^{C_\gamma}(x_1,y_1; x_2,y_2) - 8Mn < -\frac{1}{2}c_1 n^{1/3+2\delta} \right) +  4M^2n^2 \max_{0 \leq i \leq  j \leq k(n)-1}\Prob\left( L_{i,j} -8Mn > -\frac{1}{2}c_1 n^{1/3+2\delta} \right)
    \end{aligned}
\end{equation}
We again apply Proposition \ref{prop:constrained_moderate} to bound the first term. For the second term, we apply Lemma \ref{lem:diag-deficit}.
\begin{equation}
    \begin{aligned}
        &\Prob\left(\Lfull(x_1,y_1; w_{i+1},w_{i+1}) + \Lhalf(w_i,w_i;w_{j+1},w_{j+1}) + \Lfull(w_j,w_j;x_2,y_2) -8Mn > -\frac{1}{2}c_1 n^{1/3+2\delta} \right) \\
        &= \Prob\left(\Lfull(x_1,y_1; w_{i+1},w_{i+1}) - d(x_1,y_1;w_{i+1},w_{i+1}) > \frac{1}{6}c_1 n^{1/3+2\delta} \right) \\
        &+\Prob\left(\Lhalf(w_i,w_i;w_{j+1},w_{j+1}) - d(w_i,w_i;w_{j+1},w_{j+1}) > \frac{1}{6}c_1 n^{1/3+2\delta} \right) \\
        &\quad +\Prob\left(\Lfull(w_j,w_j;x_2,y_2) - d(w_j,w_j;x_2,y_2) > \frac{1}{6}c_1 n^{1/3+2\delta} \right)
    \end{aligned}
\end{equation}
Then, by Proposition \ref{prop:moderate}, Proposition \ref{prop_upper_tail}, and the cheap bound above, the above probability is bounded by $Ce^{-c {\min \{n^{2\delta}, n^{1/3} \}}}$.
\end{proof}

\section{Acknowledgments}
The author sincerely thanks their advisor, Ivan Corwin, for constant guidance and support. The author is especially grateful to Sayan Das for suggesting this problem, and to Milind Hegde for detailed feedback on an early draft of this work. The author also thanks Jiyue Zeng and Alan Zhao for valuable conversations. This research was partially supported by Ivan Corwin’s National Science Foundation grant DMS:2246576 and Simons Investigator in Mathematics award MPS-SIM-00929852.

\bibliographystyle{amsalpha}
\bibliography{bib.bib}
\end{document}
